\section{SHAP làm vũ khí}
\label{sec:ch14-shap-lam-vu-khi}

\index{tấn công né tránh|(}%
Hai mục vừa rồi đặt LIME và SHAP ở phía bị tấn công. Bài báo của mục này đảo
chiều mũi tên: ở đây SHAP nằm trong tay bên tấn công, và nạn nhân là một bộ
phân loại ảnh~\autocite{p28shapattack}.

\index{tấn công né tránh!định nghĩa}%
Loại tấn công mà bài dựng có tên riêng và cần phân biệt với hai loại kia. Một
phép \tn{tấn công né tránh}{evasion attack} sửa đầu vào ở giai đoạn suy luận,
sau khi mô hình đã huấn luyện xong, sao cho mô hình phân loại sai mà mắt người
không thấy khác biệt; nó khác phép tấn công đầu độc, vốn sửa dữ liệu huấn
luyện. Bài đặt mình ở kịch bản \tn{hộp trắng}{white-box}, tức là bên tấn công
có mô hình trong tay~\autocite{p28shapattack}. Chỗ so sánh của bài là Fast
Gradient Sign Method, viết tắt FGSM, phép tấn công né tránh quen thuộc nhất: nó
cộng vào đầu vào một lượng nhỏ theo dấu của gradient hàm mất mát, tức là đi
ngược chiều gradient descent~\autocite{p28shapattack}.

Ý của bài là thay gradient bằng giá trị SHAP. Lý do mà bài nêu là một lý do về
lượng: gradient chỉ cho biết chiều, nên FGSM nhân dấu của gradient với một hằng
số và bỏ qua độ lớn của ảnh hưởng, còn giá trị SHAP mang theo cả độ lớn, nên
bên tấn công sửa ít điểm ảnh hơn và phép sửa kín đáo
hơn~\autocite{p28shapattack}.

\index{tấn công né tránh!hình cánh bướm}%
Bước tiếp theo là quan sát thực nghiệm mà cả phép tấn công dựng trên, và nó
đáng đọc kỹ vì nó là quan sát về quan hệ giữa giá trị của một điểm ảnh và điểm
số SHAP của nó chứ không phải về bản thân điểm số. Bài tính giá trị SHAP cho
từng điểm ảnh trên hàng trăm ảnh, rồi vẽ điểm số ấy theo giá trị của chính điểm
ảnh, đã chuẩn hóa về đoạn từ 0 tới 1. Chồng nhiều ảnh lên nhau thì hiện ra một
hình mà bài gọi là hình cánh bướm: điểm số SHAP gần bằng không ở khoảng giữa
của thang giá trị và tản ra hai bên khi giá trị tiến về hai đầu mút. Với hai bộ
dữ liệu ảnh màu, vùng gần như trung hòa nằm quanh 0.4; với MNIST thì nó thấp
hơn, quanh 0.3 hoặc dưới~\autocite{p28shapattack}.

Đọc quan sát ấy theo hướng bên tấn công thì được một phép sửa: đẩy những điểm
ảnh có ảnh hưởng mạnh về phía vùng trung hòa, và như thế là rút đi thông tin mà
mô hình dựa vào. Bài phát biểu phép sửa bằng hai phương trình, và cả hai đều
cần tới điểm số SHAP của điểm ảnh đang xét, mà từ đây viết là $\phi$. Phương
trình thứ nhất đặt ra điều kiện chọn: gọi $V$ là mệnh đề $\phi_i \ge \tau$ hoặc
$\phi_i \le -\tau$, tức độ lớn của điểm số vượt một ngưỡng $\tau$. Phương trình
thứ hai nói sửa thế nào, và nó chỉ chạm tới những điểm ảnh vừa thỏa $V$ vừa nằm
sát một trong hai đầu mút của thang giá trị, cách đầu mút không quá một biên
$\eta$. Viết $x_i$ cho giá trị của điểm ảnh thứ $i$ trong ảnh $x$ và
$\tilde{x}_i$ cho giá trị của nó sau khi sửa:

\begin{equation}
\label{eq:ch14-tan-cong-shap}
\tilde{x}_i = x_i + \varepsilon \cdot
\begin{cases}
-\lvert\phi_i\rvert, & \text{nếu } x_i \ge 1-\eta \text{ và } V,\\
\phantom{-}\lvert\phi_i\rvert, & \text{nếu } x_i \le \eta \text{ và } V,
\end{cases}
\end{equation}

với $\varepsilon$ là \tn{cường độ tấn công}{attack
intensity}~\autocite{p28shapattack}. Bài viết vế trái là $x'$, chữ mà
chương~\ref{ch:lime} đã cấp cho \tn{biểu diễn có thể diễn giải}{interpretable
representation} và mà phụ lục~\ref{app:ky-hieu} cấm đổi nghĩa, nên sách viết
$\tilde{x}_i$ và giữ $x$ cho chính bức ảnh. Bài cũng viết ngưỡng là $v$ và biên
là $h$; sách đổi cả hai, vì $v(S)$ ở phụ lục~\ref{app:ky-hieu} đã thuộc về giá
trị liên minh của chương~\ref{ch:shap} và $h_x$ thuộc về ánh xạ từ biểu diễn về
đầu vào của cùng chương ấy. Ký hiệu $\varepsilon$ thì sách giữ nguyên như bài
viết, dù chương~\ref{ch:gioi-han-ly-thuyet} đã có $\varepsilon_f(b)$: cường độ
nhiễu đối kháng viết là $\varepsilon$ trong toàn bộ ngành ấy, kể cả trong
phương trình FGSM mà bài trích lại, và hai nghĩa tách nhau bằng hình dạng vì
một bên luôn có chỉ số dưới cùng đối số còn một bên luôn viết trần.

\index{tấn công né tránh!phương trình không khớp lời văn}%
Phương trình~\ref{eq:ch14-tan-cong-shap} không làm hết những gì lời văn quanh
nó mô tả, và chỗ lệch phải nói ra vì phần sau dựa vào cái phương trình làm chứ
không dựa vào cái lời văn hứa. Lời văn viết rằng phép sửa đẩy những điểm ảnh có
điểm số SHAP dương ra khỏi các vùng ảnh hưởng mạnh \emph{và} đẩy những điểm ảnh
có điểm số âm về phía hai đầu mút~\autocite{p28shapattack}. Phương trình chỉ
làm vế thứ nhất. Nhánh nào được chọn là do giá trị $x_i$ của điểm ảnh quyết
định, chứ không do dấu của $\phi_i$: một điểm ảnh sát 1 bị trừ đi và một điểm
ảnh sát 0 được cộng vào, cả hai đều dịch về giữa thang, và dấu của $\phi_i$ đã
bị dấu trị tuyệt đối xóa mất. Điều kiện $V$ cũng đối xứng theo dấu ấy, vì nó
nhận cả $\phi_i$ rất dương lẫn $\phi_i$ rất âm. Sách ghi lại cả hai, cái bài
viết và cái phương trình làm, theo đúng cách
mục~\ref{sec:ch05-integrated-gradients} đã xử lý một mệnh đề mà giả thiết của
nguồn không đỡ nổi kết luận của chính nó, và chỉ dùng phần mà phương trình đỡ
được.

\index{tấn công né tránh!kết quả}%
Bài chạy phép tấn công trên bốn bộ dữ liệu với ba kiến trúc cách xa nhau, cốt
để kết quả không lệ thuộc vào một kiến trúc. Hai bộ ảnh màu, Animal Faces và
Woman and Man Faces, dùng chung một mạng tích chập ba tầng khoảng bảy mươi sáu
nghìn tham số; Cats and Dogs Filtered dùng học chuyển giao từ EfficientNetB7
với hơn sáu mươi sáu triệu tham số; còn MNIST, bộ ảnh xám khác hẳn ba bộ kia,
dùng một mạng hai tầng tích chập~\autocite{p28shapattack}. Trong
bảng~\ref{tab:ch14-ti-le-phan-loai-sai}, chỗ cần nhìn là cặp tỉ lệ ở mỗi dòng,
và nhất là dòng duy nhất mà cặp ấy đảo chiều.

\begin{table}[htbp]
\centering
\caption{Tỉ lệ phân loại sai sau tấn công, FGSM đặt cạnh phép tấn công bằng giá
trị SHAP, cùng cường độ $\varepsilon$ mà mỗi bên cần. Hai thang $\varepsilon$
không so sánh với nhau được vì chúng nhân với hai đại lượng khác nhau. Ở ba bộ
đầu, $\varepsilon$ được chọn sao cho phép sửa còn khó thấy bằng mắt; ở MNIST
thì bài ghi rằng không giấu được, nên hai bên chỉ được chỉnh cho mức lộ ngang
nhau. Cột cuối là độ chính xác của chính mô hình trước khi bị tấn công, để thấy
cả bốn mô hình đều là mô hình tốt. Số liệu từ bảng 1 của bài
báo~\autocite{p28shapattack}.}
\label{tab:ch14-ti-le-phan-loai-sai}
\begin{tabular}{@{}lccccc@{}}
\toprule
Bộ dữ liệu & FGSM & $\varepsilon$ & Tấn công SHAP & $\varepsilon$ & Độ chính xác mô hình \\
\midrule
Animal Faces & 52\% & 0.2 & 73\% & 50 & 99.5\% \\
Cats and Dogs Filtered & 98\% & 0.01 & 89\% & 80 & 98\% \\
MNIST & 34\% & 0.2 & 60\% & 1.5 & 99\% \\
Woman and Man Faces & 69\% & 0.2 & 98\% & 60 & 97\% \\
\bottomrule
\end{tabular}
\end{table}

Phép tấn công bằng SHAP thắng ở ba bộ trên bốn và thua ở bộ còn lại, Cats and
Dogs Filtered, với 89 phần trăm so với 98 phần trăm. Cho nên câu đọc được từ
bảng không phải câu rằng phép tấn công mới luôn mạnh hơn. Câu mà bài phát biểu
hẹp hơn thế và bảng đỡ được: tỉ lệ của FGSM trải từ 34 tới 98 phần trăm còn tỉ
lệ của phép tấn công SHAP trải từ 60 tới 98, nên phép tấn công bằng SHAP cho
kết quả ổn định hơn~\autocite{p28shapattack}. Phần tóm tắt của bài nói rõ điều
kiện mà nó nhắm tới, là những tình huống gradient bị
che~\autocite{p28shapattack}: khi gradient triệt tiêu, bùng nổ hoặc tản mát thì
hướng mà FGSM đọc được không còn mang thông tin, trong khi giá trị SHAP không
đọc hàm mất mát nên không chịu ảnh hưởng ấy.

\index{tấn công né tránh!là một phép đo về attribution}%
Còn một cách đọc bảng~\ref{tab:ch14-ti-le-phan-loai-sai} nữa mà bài không đưa
ra, và nó là quan sát của sách. Chữ \enquote{faithful} không xuất hiện lần nào
trong bài, và bài chỉ nhắc tới attribution đặc trưng đúng một lần, trong phần
nền. Nhưng phép tấn công chỉ chạy được nếu giá trị SHAP chỉ đúng vào những điểm
ảnh mà mô hình thật sự dựa vào: sửa nhầm chỗ thì mô hình không đổi kết luận. Tỉ
lệ phân loại sai vì thế cũng là một phép đo về chất lượng của chính bộ điểm số
đã dẫn đường cho phép sửa.

Đại lượng ấy có mấy tính chất mà không họ chỉ số nào của
chương~\ref{ch:do-do-trung-thuc} có đủ cùng lúc. Nó không cần ground truth, vì
mô hình tự chấm điểm bằng đầu ra của chính nó. Nó không cần thang quy ước, vì
phân loại sai hay không là một sự kiện quan sát được. Và nó không tự chọn phép
nhiễu cho có lợi, vì phép nhiễu ở đây phải đạt một mục tiêu độc lập là đổi được
kết luận của mô hình. Về hình dạng thì nó chính là họ xóa dần của
mục~\ref{sec:ch08-ho-xoa-dan} chạy dưới dạng đối kháng: cùng một vòng lặp là bỏ
đi thông tin ở những chỗ điểm cao rồi xem mô hình đổi bao nhiêu, chỉ khác ở chỗ
người chạy có mục tiêu và có ngân sách nhiễu.

\index{mức quan trọng!đo được qua tấn công}%
Chỗ dừng của cách đọc ấy cũng phải nói ra trong cùng một hơi, vì bỏ qua nó là
mắc đúng lỗi mà bài báo neo chẩn đoán. Cái mà tỉ lệ phân loại sai chứng thực là
\tn{mức quan trọng}{importance}, tức tính chất ảnh hưởng nhân quả mà
mục~\ref{sec:ch09-chan-doan} đã tách riêng ra và giữ cách khỏi độ trung thực:
sửa những điểm ảnh này thì đầu ra đổi, cho nên điểm số có bắt được ảnh hưởng.
Nó không nói rằng bộ điểm số phản ánh đúng quá trình mô hình dùng để tới đầu
ra, và với một mạng tích chập thì hai điều ấy cách nhau khá xa. Gộp chúng lại
là đúng phép gộp mà bài báo neo của
chương~\ref{ch:chi-so-khong-do-do-trung-thuc} chẩn đoán trên \tn{chuỗi suy
luận}{chain of thought}, nên sách để phép đo này ở đúng chỗ của nó: một điều
kiện cần đo được rẻ, không phải một phép đo độ trung thực.

\index{tấn công né tránh|)}%
